\chapter{Linear Logic, Cut Elimination, and Zero Allocation}\label{ch:linear}

Linear logic is the logic of resources used exactly once \citep{girard}: its formulas cannot be duplicated or discarded freely. A dataplane can promise zero allocation on the hot path when every datum is produced once and consumed once, buffers are preallocated, and intermediate values of a pipeline are held in registers rather than heap buffers. This chapter proves that the composite Mealy machine of \autoref{def:comp} is a fused pipeline whose intermediate value is a register (\autoref{thm:49}), that linear atoms together with preallocated rings force zero hot-path allocation (\autoref{thm:50}), and that a linear combinator calculus of molecules, matching the verifier \texttt{affine\_typer}, never duplicates subterms and never grows the syntax tree (\autoref{thm:55}). The Girard and Wadler readings are analogies: product-state composition is the operational counterpart of deforestation for Mealy machines, not a transcription of cut elimination in proof nets \citep{girard, wadler}.

\paragraph{Pipelines and fusion.}
A pipeline is a composite $p = g_n \circ \cdots \circ g_1 : A_0 \to A_n$ of molecules $g_i : A_{i-1} \to A_i$. By \autoref{def:comp} the composite is a single machine whose state space is the product of the component state spaces and whose transition runs the components in order, feeding each output into the next input inside one step. The intermediate value of type $A_i$ exists as a register for the duration of that step and is never stored as a heap buffer in this realization.

\setcatalog{49}
\begin{theorem}[Product-state fusion (\autoref{thm:49})]\label{thm:49}
Every pure pipeline $p = g_n \circ \cdots \circ g_1 : A_0 \to A_n$ in Mol is equal to a single molecule $h : A_0 \to A_n$, the \emph{fused pipeline}, with the following properties:
\begin{enumerate}
\item $h \simeq p$: fusion preserves behavior;
\item the intermediate types $A_1, \ldots, A_{n-1}$ never materialize as buffers during the execution of $h$: no datum of type $A_i$ is stored between stages;
\item $h$ is constructed from the presentations of the $g_i$ by an explicit finite procedure of worst-case complexity $O\bigl (|S_1| \cdots |S_n| \cdot (|A_0| + \cdots + |A_{n-1}|)\bigr)$, where $|S_i|$ is the size of the state space of $g_i$.
\end{enumerate}
\end{theorem}

\begin{proof}
First fuse two molecules. Let $f : A \to B$ and $g : B \to C$ be presented by $(S_f, s_f, \delta_f)$ and $(S_g, s_g, \delta_g)$. Define $h : A \to C$ with state space $S_f \times S_g$, initial state $(s_f, s_g)$, and transition
\[
\delta_h\bigl ((u, v), a\bigr) = \bigl (c, (u', v')\bigr),
\qquad \text{where } (b, u') = \delta_f (u, a), \quad (c, v') = \delta_g (v, b) .
\]
This is a valid Mealy machine, and by the definition of composition in Mol (\autoref{ch:mol}) it is exactly the composite $g \circ f$: on every input word, $h$ follows the same state trajectory $(u, v) \mapsto (u', v')$ and emits the same output word as the two-stage machine, so $h \simeq g \circ f$. In a single transition of $h$, the value $b$ is produced by $\delta_f$ and consumed by $\delta_g$ within that same transition; it is held in a register and never stored, so no buffer of type $B$ exists in this realization.

Now iterate. By induction on $n$, the pipeline $g_{n-1} \circ \cdots \circ g_1$ is equal to a single molecule $h' : A_0 \to A_{n-1}$ with state space $S_1 \times \cdots \times S_{n-1}$ and no intermediate buffers; fusing $h'$ with $g_n$ by the two-molecule construction gives the single molecule $h$ with state space $S_1 \times \cdots \times S_n$, behavior equal to $p$, and no buffer of any type $A_i$ ($1 \leq i \leq n-1$). The transition table of $h$ has $|S_1| \cdots |S_n| \cdot |A_0|$ entries, each computed by one step of every $g_i$ (work $O (|A_0| + \cdots + |A_{n-1}|)$), giving the stated construction cost. The construction is the product-state composition of \autoref{def:comp}, iterated. In programming terms it is deforestation of the intermediate buffer \citep{wadler}; it is the operational analogue, for Mealy machines, of composing morphisms without materializing the intermediate object, and is not identified with Girard's cut-elimination theorem for proof nets \citep{girard}.
\end{proof}

\paragraph{Remark (fusion cost).}
The fused machine is a morphism of Mol, so by the additivity of the cost model (\autoref{thm:25}) its cost is at most the sum of the component costs; because the intermediate value is transient, the buffer copy that a two-stage realization pays is absent, which is the content of the fusion bound (\autoref{thm:27}): $c (g \circ f) \leq c (f) + c (g)$, strict when the intermediate type never materializes. \autoref{thm:49} supplies the mechanism behind that bound.

\paragraph{Linearity and ownership.}
Call an atom $f : A \to B$ \emph{linear} if for every state $s$ and every input $a$ in its domain, the transition $(b, s') = \delta (s, a)$ reads $a$ exactly once and produces exactly one output $b$. A linear atom never copies its input, never drops it, and never holds it past the transition that consumes it; at every instant each datum is \emph{owned} by exactly one atom. This is the ownership discipline of Rust restated as a property of molecules.

\setcatalog{50}
\begin{theorem}[Linearity forces zero allocation (\autoref{thm:50})]\label{thm:50}
Let $M$ be a molecule built from linear atoms by the operations $\circ$, $\otimes$, and the batch construction of \autoref{ch:operad}. Then:
\begin{enumerate}
\item $M$ is linear on each input factor: a morphism $A \to B$ consumes every component of an $A$-value exactly once and produces every component of a $B$-value exactly once;
\item along every execution of a sequential pipeline of linear atoms, each live token is owned by exactly one stage; a tensor $f \otimes g$ on a pair consumes two tokens and produces two tokens, so the token count equals the arity of the live type, not the constant $1$;
\item consequently, if the buffers of $M$ are preallocated to their maximum occupancy (bounded by the ring invariant, \autoref{thm:48}), the hot path performs no allocation and no deallocation.
\end{enumerate}
\end{theorem}

\begin{proof}
(i) By induction on the construction of $M$. An atom is linear by hypothesis. Composite: in $g \circ f$, an input token $a$ of type $A$ is read exactly once by $f$, producing exactly one token $b$ of type $B$; $b$ is read exactly once by $g$, producing exactly one token $c$ of type $C$. Tensor: $f \otimes g$ on a pair $(a, a')$ reads $a$ exactly once in $f$ and $a'$ exactly once in $g$ and emits exactly one pair. Batch: $f^{(n)}$ reads each of its $n$ inputs exactly once by the definition of the batch transition and emits exactly $n$ outputs. Every molecule is obtained by finitely many such steps, so (i) holds.

(ii) Fix a sequential pipeline execution and let $L (t)$ be the number of live data tokens in the preallocated rings at time $t$, counted by the arity of the live type. A sequential step of a linear atom $A \to B$ consumes the $A$-arity and produces the $B$-arity; tokens move from the input ring to the output ring and are never copied. A tensor step consumes and produces a pair, so $L$ equals $2$ while a pair is live. In all cases $L (t)$ is bounded by the sum of the arities of the live types, independently of the length of the run.

(iii) By (ii) occupancy never exceeds $L (0)$. If each stage buffer is preallocated with capacity at least its share of $L (0)$ (a ring of capacity $2^k \geq L (0) + 1$ per stage satisfies the in-flight bound $w - r \leq 2^k - 1$ of \autoref{thm:48}), then every write fits in preallocated storage and every read empties a slot in place. No token is copied, created, or destroyed, so the execution acquires and releases no heap memory: zero allocations and zero deallocations on the hot path.
\end{proof}

\paragraph{Geometry of interaction.}
In Girard's geometry of interaction a proof is interpreted as an operator, and cut elimination is the \emph{execution} of that operator: a token travels through the proof, and cuts are crossed by iterating the operator until the token exits \citep{girard, arxiv-linear-logic-goi}. The reading in Mol is literal: a datum is a token, an atom is a transition $\delta$ that carries the token from its input port to its output port, and a molecule is the composition of these port-to-port maps. The fusion theorem (\autoref{thm:49}) is the execution: it computes in one pass the effect of the whole pipeline on any token, without materializing the cut formulas. The token interpretation is also the performance claim: the working set is the set of live tokens, invariant by \autoref{thm:50}, which bounds the cache-miss component of the cost vector of \autoref{ch:enrich}.

\paragraph{Curry--Howard for molecules.}
Read propositions as types and proofs as molecules: the linear implication $A \multimap B$ is the type of molecules $A \to B$, and composition is the cut rule; a well-typed dataplane element is a proof of its interface type. The theorems of this chapter upgrade this typing discipline to a resource guarantee: if the type system records that every atom is linear (the hypothesis of \autoref{thm:50}) and that fusion removes intermediate structure (\autoref{thm:49}), then a well-typed hot path is allocation-free by construction. This is the categorical content of ownership in Rust: a discipline of exact-once consumption is precisely the data that parts (ii)--(iii) of \autoref{thm:50} require, so the guarantee is enforced at compile time rather than by runtime measurement \citep{wadler}.

\section{The Linear Combinator Calculus}\label{sec:linear-affine}

The ownership discipline of \autoref{thm:50} is a semantic hypothesis on atoms. This section enforces a corresponding discipline on \emph{combinator terms}: a linear calculus $\Lambda$ of molecules, with no weakening and no contraction, whose sequential composition is pipeline composition $t ; u = u \circ t$ and whose tensor is the bifunctor of \autoref{thm:3}. Affine logic would permit weakening (dropping a resource); $\Lambda$ forbids it, so the system is linear. The verifier \texttt{affine\_typer} of \autoref{app:a} implements this calculus (the binary keeps its historical name). Reduction never duplicates a subterm and never increases the number of syntax-tree nodes. That bound is a bound on the combinator term, not on the dataplane heap; zero hot-path heap remains \autoref{thm:50}, which requires preallocated rings.

\begin{definition}[Linear combinator calculus]\label{def:lambda}
The types of $\Lambda$ are generated from a fixed set $\mathsf{B}$ of \emph{base types} by the grammar
\[
A, B \;\;::=\;\; \beta \;\mid\; A \multimap B \;\mid\; A \otimes B \;\mid\; I,
\qquad \beta \in \mathsf{B},
\]
where $A \multimap B$ is the type of molecules $A \to B$, $A \otimes B$ the tensor of types, and $I$ the unit. A \emph{linear context} is a finite partial map from variables to types; contexts are combined by the disjoint union $\Gamma, \Delta$, defined only when the domains of $\Gamma$ and $\Delta$ are disjoint. The \emph{terms} are generated by
\[
t, u \;\;::=\;\; x \;\mid\; a \;\mid\; t ; u \;\mid\; t \otimes u \;\mid\; (),
\]
where $x$ ranges over variables standing for assumed molecules and $a$ ranges over a fixed signature $\Sigma$ of atom constants $a : A \multimap B$. Sequential composition is in pipeline order: $t ; u$ means ``run $t$, then $u$'', the composite $u \circ t$ of Mol. Tensor of terms is the bifunctor. The typing rules are:
\begin{enumerate}
\item[(var)] $x : A \vdash x : A$;
\item[(atom)] $\vdash a : A \multimap B$ for every atom $a : A \multimap B$ in $\Sigma$;
\item[(comp)] from $\Gamma \vdash t : A \multimap B$ and $\Delta \vdash u : B \multimap C$ with $\Gamma$ and $\Delta$ disjoint, $\Gamma, \Delta \vdash t ; u : A \multimap C$;
\item[(tensor)] from $\Gamma \vdash t : A \multimap B$ and $\Delta \vdash u : C \multimap D$ with $\Gamma$ and $\Delta$ disjoint, $\Gamma, \Delta \vdash t \otimes u : (A \otimes C) \multimap (B \otimes D)$;
\item[(unit)] $\vdash () : I \multimap I$.
\end{enumerate}
There are no rules of weakening or contraction. A \emph{molecule term} is a closed well-typed term $\vdash t : A \multimap B$; these are the morphisms of Mol in the fragment generated by composition and tensor.
\end{definition}

\begin{definition}[Operational semantics of $\Lambda$]\label{def:lambda-op}
The reduction relation $t \longrightarrow t'$ on well-typed terms is the smallest relation closed under the term constructors (if $t \longrightarrow t'$ then $t ; u \longrightarrow t' ; u$, $u ; t \longrightarrow u ; t'$, $t \otimes u \longrightarrow t' \otimes u$, and $u \otimes t \longrightarrow u \otimes t'$) and containing the following redex rules.
\begin{enumerate}
\item[(CE)] \emph{Composition elimination} (associativity of the cut): $(t ; u) ; v \longrightarrow t ; (u ; v)$.
\item[(TE)] \emph{Tensor elimination} (interchange): for $t : A \multimap B$, $u : C \multimap D$, $v : B \multimap E$, $w : D \multimap F$,
\[
(t \otimes u) ; (v \otimes w) \longrightarrow (t ; v) \otimes (u ; w).
\]
\item[(UE)] \emph{Unit elimination}: $t \otimes () \longrightarrow t$ and $() \otimes t \longrightarrow t$, using $A \otimes I \equiv A \equiv I \otimes A$.
\end{enumerate}
Rule (CE) reassociates a pipeline; rule (TE) is the interchange law of \autoref{thm:11}; rule (UE) tensors with the identity on $I$. Every rule consumes exactly one redex and its right-hand side is built from subterms of that redex.
\end{definition}

\begin{lemma}[Subject reduction]\label{lem:subject-reduction}
If $\Gamma \vdash t : A$ and $t \longrightarrow t'$, then $\Gamma \vdash t' : A$.
\end{lemma}

\begin{proof}
It suffices to check the redex rules; the congruence steps then follow by the typing rules. Rule (CE): if $(t ; u) ; v : A \multimap D$, then $t : A \multimap B$, $u : B \multimap C$, and $v : C \multimap D$; the reduct $t ; (u ; v)$ derives $A \multimap D$ by two applications of (comp). Rule (TE): with the displayed types, the left-hand side has type $(A \otimes C) \multimap (E \otimes F)$; the reduct $(t ; v) \otimes (u ; w)$ has components of type $A \multimap E$ and $C \multimap F$ by (comp), hence the same type by (tensor). Rule (UE): $t \otimes ()$ has type $(A \otimes I) \multimap (B \otimes I)$, identified with $A \multimap B$, and the reduct $t$ has type $A \multimap B$. Every redex rule therefore preserves the judgment.
\end{proof}

\setcatalog{55}
\begin{theorem}[The no-allocation theorem, \autoref{thm:55}]\label{thm:55}
Let $\Lambda$ be the linear combinator calculus of \autoref{def:lambda} with the operational semantics of \autoref{def:lambda-op}. Then:
\begin{enumerate}
\item[(i)] \emph{Linearity invariant.} In every derivation of a judgment $\Gamma \vdash t : A$, every variable of $\Gamma$ occurs exactly once in $t$.
\item[(ii)] \emph{No duplication.} Every reduction rule is non-duplicating: no rule copies a subterm.
\item[(iii)] \emph{No fresh syntax-tree nodes.} The evaluation of every well-typed closed molecule term is a finite reduction sequence, and every step consumes one redex and produces only subterms already present in the redex. The node set of every reduct is a subset of the node set of the initial term, so the combinator term never grows. This is a bound on the syntax tree of $\Lambda$, verified by \texttt{affine\_typer}. It does not by itself bound the dataplane heap: zero hot-path heap is \autoref{thm:50}, which requires linear atoms and preallocated rings.
\item[(iv)] Exact-once use of combinator variables is enforced by typing. Rust's ownership (move semantics as linear consumption) is the operational realization of the same discipline on values, which is the hypothesis of \autoref{thm:50}.
\end{enumerate}
\end{theorem}

\begin{proof}
(i) By induction on the derivation of $\Gamma \vdash t : A$. Rule (var): $x$ occurs exactly once. Rules (atom) and (unit): the context is empty. Rule (comp): $\Gamma, \Delta \vdash t ; u : A \multimap C$ with disjoint domains; by induction every variable of $\Gamma$ occurs exactly once in $t$ and every variable of $\Delta$ exactly once in $u$, hence exactly once in $t ; u$. Rule (tensor) is the same argument. The system therefore admits neither contraction nor weakening.
(ii) Rule (CE): the redex $(t ; u) ; v$ contains $t$, $u$, $v$ once each, and the reduct $t ; (u ; v)$ contains exactly those subterms. Rule (TE): the redex $(t \otimes u) ; (v \otimes w)$ contains $t, u, v, w$ once each; the reduct contains the same four subterms. Rule (UE): the reduct is $t$, and the erased unit carries no variable. No rule copies a subterm. By subject reduction every reduct is well-typed, so a duplicating step would contradict (i).
(iii) Termination: the lexicographic measure $M (t) = (q (t), p (t), |t|)$ of the original argument is unchanged in its reasoning. Rule (CE) decreases $q$ or else $p$; rule (TE) decreases $q$ by $1$; rule (UE) decreases $|t|$. The triple lies in the well-founded set $\mathbb{N}^{3}$, so every reduction sequence is finite. Fresh nodes: each redex rule reuses the subterms of the redex and introduces no constructor that was absent from the redex. By induction along the sequence, every node of every reduct is a node of the initial term. The peak syntax-tree size is therefore $|t_0|$. The claim is about combinator terms. Dataplane heap allocation is excluded only when \autoref{thm:50} applies: linear atoms and rings preallocated to the occupancy bound of \autoref{thm:48}.
\end{proof}

\paragraph{Ownership in Rust.}
\autoref{thm:50} assumes exact-once consumption of values; \autoref{thm:55} proves exact-once use of combinator variables and a non-growing syntax tree. Rust's move semantics realize the value-level discipline. The two theorems together, not \autoref{thm:55} alone, are the zero-allocation claim of the hot path.

\paragraph{Remark (scope of the guarantee).}
\autoref{thm:55} is a guarantee about the fragment of Mol generated by composition, tensor, and the unit: the linear combinator core. The buffer-level guarantee remains \autoref{thm:50}.
